\documentclass[12pt]{article}
\usepackage[T1]{fontenc}
\usepackage[utf8]{inputenc}
\usepackage{lmodern}
\usepackage{textcomp}
\usepackage{enumitem}
\usepackage{geometry}
\usepackage{graphicx}
\usepackage{amsmath, amssymb}
\geometry{margin=1in}
\begin{document}
\begin{center}
Economics - Graduate Level Assessment 19 \
Total Marks: 40 \
Answer all questions.
\end{center}

\section*{Multiple Choice (10 marks)}
\begin{enumerate}[label=\arabic*.]
\item What is meant by near-money? Give an example.
\begin{enumerate}[label=(\alph*)]
    \item Near-monies are physical currencies
    \item Near-monies are long-term investment bonds
    \item Near-monies are commodities used for barter
    \item Near-monies are precious metals such as gold and silver
    \item Near-monies are ownership stakes in real estate
\end{enumerate}
\item Assume the reserve requirement is five percent. If the FED sells \$10 million worth of government securities in an open market operation then the money supply can potentially
\begin{enumerate}[label=(\alph*)]
    \item decrease by \$200 million.
    \item decrease by \$100 million.
    \item increase by \$200 million.
    \item decrease by \$10 million.
    \item increase by \$500 million.
\end{enumerate}
\item Which of the following is true of money and financial markets?
\begin{enumerate}[label=(\alph*)]
    \item As the demand for bonds increases the interest rate increases.
    \item For a given money supply if nominal GDP increases the velocity of money decreases.
    \item When demand for stocks and bonds increases the asset demand for money falls.
    \item A macroeconomic recession increases the demand for loanable funds.
\end{enumerate}
\item Compared to perfect competition in the long run, monopoly has
\begin{enumerate}[label=(\alph*)]
    \item more price stability.
    \item more choices of products for consumers.
    \item lower prices.
    \item more efficiency.
    \item less market control.
\end{enumerate}
\item Which of the following is equal to one?
\begin{enumerate}[label=(\alph*)]
    \item The elasticity of the long-run aggregate supply curve
    \item The spending (or expenditure) multiplier
    \item The investment multiplier
    \item The balanced budget multiplier
    \item The aggregate demand multiplier
\end{enumerate}
\end{enumerate}

\section*{True / False (10 marks)}
\textit{Answer True or False}
\begin{enumerate}[label=\arabic*.]
\setcounter{enumi}{5}
\item True or False: The Rule of 72 is a banking regulation that dictates how interest rates should be applied to savings accounts.
\item True or False: When the value of the U.S. dollar appreciates, U.S. residents are likely to take more vacations in foreign countries due to lower relative costs.
\item True or False: When GNP is $2,050,000 and capital consumption is $500,000, the NNP can be calculated as \$1,950,000.
\item True or False: Raising taxes on disposable income is an effective measure to combat built-in inflation by reducing consumer spending.
\item True or False: An expectation of rising prices for a product typically leads to an increase in current demand for that product.
\end{enumerate}

\section*{Long Form Response (20 marks)}
\textit{Answer each question with a clear and complete explanation.}
\begin{enumerate}[label=\arabic*.]
\setcounter{enumi}{10}
\item Discuss how a higher rate of inflation relative to other nations can impact a nation's trade deficit, including the mechanisms involved and potential long-term effects.
\item Discuss how increasing costs, indicative of decreasing returns to scale, influence the degree of specialization within an economy. Provide examples to support your analysis.
\end{enumerate}

\vfill
\noindent\textit{End of Paper}
\end{document}